\documentclass{article}
\usepackage{amsmath}
\usepackage{amsfonts}
\usepackage{amssymb}

\title{The Cone Generated by Positive Semidefinite Mean Polynomials}
\author{Mehdi Ghasemi, Salma Kuhlmann}
\date{\today}

\usepackage{amsthm}
\newtheorem{theorem}{Theorem}
\newtheorem{proposition}{Proposition}
\newtheorem{remark}{Remark}
\usepackage{xcolor}
\usepackage{hyperref}

\begin{document}
\maketitle

\begin{abstract}
Certain classes of (homogeneous) polynomials, such as Positive Semidefinite (PSD) Diagonal minus Tail (D-T)
forms and those derived from the Arithmetic-Geometric Mean inequality, like Sums Of Nonnegative Circuit (SONC)
polynomials, can be identified with nonnegative mean polynomials. This paper studies PSD forms that are weighted
power means, denoted by $M_{q,p}(Y, w)$, and investigates the cone of nonnegative mean real $n$-variate forms
of degree $2d$, denoted $M_{n,2d}(X, \alpha)$, and its relation to other important cones: the cone of all PSD
forms, the cone of Sums of Squares (SOS), and the cone of SONC polynomials. A key object of study is the sum
of the SOS and SONC cones, known as the SOSONC cone.
\end{abstract}

\section{Introduction}\label{s00}

Certain classes of (homogeneous) polynomials, such as Positive Semidefinite (PSD) Diagonal minus Tail (D-T) forms and
those derived from the Arithmetic-Geometric Mean inequality, like Sums Of Nonnegative Circuit (SONC) polynomials, can
be identified with nonnegative mean polynomials.

In this paper, we study PSD forms that are weighted power means, denoted $M_{q,p}(Y, w)$, where $Y = (Y_1, \dots, Y_m)$
is an $m$-tuple of symbols and $w = (w_1, \dots, w_m) \in \mathbb{R}_+^m$ is a vector of positive weights. In particular,
we investigate the cone of nonnegative mean real $n$-variate forms of degree $2d$, denoted $M_{n,2d}(X, \alpha)$, where
$X = (X_1, \dots, X_n)$ and $\alpha \in \mathbb{N}_0^n$.

We study the relationships between the following cones:
\begin{itemize}
   \item $\mathcal{P}_{n,2d}$: the cone of all positive semidefinite (PSD) forms in $n$ variables of degree $2d$; our cone
   is a subset: $\mathcal{M}_{n,2d} \subseteq \mathcal{P}_{n,2d}$.
   \item $\Sigma_{n,2d}$: the cone of sums of squares (SOS) polynomials.
   \item $\mathcal{C}_{n,2d}$: the cone of sums of nonnegative circuit (SONC) polynomials.
\end{itemize}

A key object of study is the sum of the SOS and SONC cones, known as the \emph{SOSONC cone}, defined by M. Schick:
\[
   SOSONC = (\Sigma + \mathcal{C})_{n,2d}
\]

In \S\ref{s01} (Preliminaries on mean polynomials), we introduce the preliminary concepts of weighted power means,
denoted $M_{q,p}(Y,w)$, $q, p \in \mathbb{N}_0$, $Y = (Y_1, \dots, Y_m)$, $w = (w_1, \dots, w_m) \in \mathbb{R}_+^m$, and
the fundamental tool of Jensen's inequality. We prove (see Proposition 1.1) that if $p < q$, then $M_{q,p}(Y,w)$ is PSD.
We identify certain classes of nonnegative polynomials as instances of $M_{q,p}(Y,w)$ for appropriate choices of $q, p, Y$,
and $w$.

In \S\ref{s02} (Nonnegative Circuit Polynomials), we re-prove \cite[Corolary 2.2]{GM01} and \cite[Theorem 2.4]{DIdW01}
using mean polynomials, giving necessary and sufficient conditions for nonnegativity of D-T forms and circuit polynomials.

In \S\ref{s03} (The cones $M_{2d,0}(\underline{X},\alpha)$ and $M_{1,0}(Y,\lambda)$), we prove Proposition 3.1: the
forms $M_{2d,0}(X,\alpha)$ are PSD D-T for any $d \in \mathbb{N}$ and any $\alpha \in \mathbb{N}_0^n$ with $|\alpha| = 2d$.
The cone generated by these forms is the cone of sums of PSD D-T forms.

In Theorem 3.1, we prove that the SONC cone, $\mathcal{C}_{n,2d}$, is the cone generated by the forms $M_{1,0}(Y,\lambda)$.

In \S\ref{s04}, we study $M_{2d,p}(X,\alpha)$ with $|\alpha| = 2d$, $p | 2d$ in connection to SOBS and SOS. We prove
in Proposition 4.1 that $M_{2d,0}(X,\alpha)$ is SOBS and PSD D-T, $|\alpha| = 2d$. In Theorem 4.1, we continue investigating
$M_{2d,p}(X,\alpha)$ with $d = 1, 2, 3$, $p | 2d$ and prove that they are SOS for all $n \in \mathbb{N}$ (by reduction to
Hilbert's case and induction on the number of variables).

In \S\ref{s05} (General Case: Beyond the SOSONC cone), the ultimate goal is to move towards a comprehensive
understanding of the general form $M_{q,p}(Y,w)$. In particular, we examine the Choi-Lam form $Q$ and the second Robinson
form $\hat{R}$ (see Section 5.1).

In \cite{MSc01}, it is observed that $Q$ is SONC \cite[page 31]{MSc01}, but is not SOSONC \cite[Lemma 4.2.5, page 107]{MSc01}.
It was further observed by Reznick \cite{Rez07} that $\hat{R}$ can be obtained from $Q$ by a linear change of variables.
In Section 5.2, we establish moreover that
\[
   Q(x, y, z, w) = 4 M_{1,0}((x^4, y^4, z^4, w^4), (1, 1, 1, 1))
\]
and accordingly
$$\hat{R}(x,y,z,w)=8M_{1,0}(G, (1,1,1,1))$$
where $G=\frac{1}{2}(x+w-y-z, y+w-x-z, z+w-x-y, w-(x+y+z))$. Mehdi to check the computation.
Therefore $\hat{R}$ is a separating form for the cone of mean polynomials and the SOSONC cone.
In \S\ref{s06} (Positivestellensatz and relation the the cone of $2d$ powers).

In \S\ref{s07} (Application to Polynomial Optimization): Mehdi: to check the algorithmic aspect of means



\section{Preliminaries on Mean Polynomials}\label{s01}

Let $\varphi(x) = x^{\frac{2q}{p}}$, where $q > p$. The second derivative of $\varphi$ is
\[
   \varphi''(x) = \frac{q}{p}\left(\frac{q}{p}-1\right)x^{\frac{2q}{p}-2}
\]
Since $q > p$, the coefficient $\frac{q}{p}\left(\frac{q}{p}-1\right)$ is always positive, and if $q$ is even then
$\varphi''$ is always positive, so $\varphi$ is convex. By Jensen's inequality,
\[
   \varphi\left(\frac{\sum w_i x_i}{\sum w_i}\right) \leq \frac{\sum w_i \varphi(x_i)}{\sum w_i}
\]
where $w_i \in \mathbb{R}_+$. Replacing $x_i$ with $x_i^p$, we get
\[
   \varphi\left(\frac{\sum w_i x_i^p}{\sum w_i}\right) \leq \frac{\sum w_i \varphi(x_i^p)}{\sum w_i}
\]
and hence
\[
   \left(\frac{\sum w_i x_i^p}{\sum w_i}\right)^{\frac{q}{p}} \leq \frac{\sum w_i x_i^{q}}{\sum w_i}
\]
or equivalently $M_p(x, w) \leq M_q(x, w)$. Raising both sides to the power $c = \mathrm{lcm}(q, p)$, we conclude that
the function
\[
   \mathcal{M}_{q,p} = M_q^c(x, w) - M_p^c(x, w)
\]
is always a nonnegative polynomial function in $x = (x_1, \dots, x_n)$.

\subsection*{Basic properties of $M_p$}
\begin{itemize}
   \item $\displaystyle \lim_{p \rightarrow -\infty} M_p(\mathbf{x}, w) = \min(x_1, \dots, x_n)$.
   \item $\displaystyle \lim_{p \rightarrow \infty} M_p(\mathbf{x}, w) = \max(x_1, \dots, x_n)$.
   \item $\displaystyle \lim_{p \rightarrow 0} M_p(\mathbf{x}, w) = \left( \prod_{i=1}^n x_i^{w_i} \right)^{1/\sum_{i=1}^n w_i}$.
   \item If $p < q$, then $M_p(\mathbf{x}, w) \leq M_q(\mathbf{x}, w)$. Equality holds when $x_1 = x_2 = \dots = x_n$.
\end{itemize}


\section{Nonnegative Diagonal minus Tail and Circuit Polynomials}\label{s02}
~~
\begin{theorem}
Let the function $f(x)$ be defined as
\[
   f(x) = \sum_{i=1}^n \beta_i x_i^{2d} - \mu x^{\alpha}
\]
where $x = (x_1, \dots, x_n)$ with $x_i \geq 0$, $\beta_i > 0$, and $\alpha = (\alpha_1, \dots, \alpha_n)$ with
$\alpha_i > 0$. The monomial $x^{\alpha}$ is defined as $x^{\alpha} = \prod_{i=1}^n x_i^{\alpha_i}$. We assume the
degree is conserved, i.e., $\sum_{i=1}^n \alpha_i = 2d$.

The function $f(x)$ is positive semidefinite (PSD), meaning $f(x) \geq 0$ for all $x \in \mathbb{R}_+^n$, if and only
if the following condition holds:
\[
   \mu^{2d} \prod_{i=1}^n \alpha_i^{\alpha_i} \leq (2d)^{2d} \prod_{i=1}^n \beta_i^{\alpha_i}
\]
\end{theorem}

\begin{proof}
The proof consists of two parts for the "if and only if" statement.

	\textbf{(\(\Leftarrow\))} Assume the inequality above holds. We want to show that $f(x) \geq 0$. This is equivalent to showing
\[
   \sum_{i=1}^n \beta_i x_i^{2d} \geq \mu x^{\alpha} \quad \forall x \in \mathbb{R}_+^n
\]
The proof relies on the weighted Arithmetic Mean-Geometric Mean (AM-GM) inequality. Applying this to the terms
$y_i = \frac{\alpha_i}{\beta_i} x_i^{2d}$ with weights $w_i = \alpha_i$ (where $\sum w_i = 2d$), we get
\[
   \frac{\sum_{i=1}^n \beta_i x_i^{2d}}{2d} \geq \left( \prod_{i=1}^n \left( \frac{\alpha_i}{\beta_i} \right)^{\alpha_i} \right)^{1/2d} x^{2d}
\]
which implies
\[
   \sum_{i=1}^n \beta_i x_i^{2d} \geq 2d \left( \prod_{i=1}^n \left( \frac{\alpha_i}{\beta_i} \right)^{\alpha_i} \right)^{1/2d} x^{\alpha}
\]
From our initial assumption, we can take the $2d$-th root and rearrange to find
\[
   \mu \leq 2d \left( \prod_{i=1}^n \left( \frac{\alpha_i}{\beta_i} \right)^{\alpha_i} \right)^{1/2d}
\]
Combining these results, we chain the inequalities
\[
   \mu x^{\alpha} \leq 2d \left( \prod_{i=1}^n \left( \frac{\beta_i}{\alpha_i} \right)^{\alpha_i} \right)^{1/2d} x^{\alpha} \leq \sum_{i=1}^n \beta_i x_i^{2d}
\]
Thus, we have shown $\sum_{i=1}^n \beta_i x_i^{2d} \geq \mu x^{\alpha}$, which holds for $x_i > 0$. If any $x_i = 0$,
then $x^{\alpha} = 0$ and the inequality becomes $\sum \beta_i x_i^{2d} \geq 0$, which is trivially true.

	\textbf{(\(\Rightarrow\))} Now, we prove the other direction by contradiction. Assume that $f(x)$ is PSD, but the
   inequality above does \emph{not} hold. This means
\[
   \mu^{2d} \prod_{i=1}^n \alpha_i^{\alpha_i} > (2d)^{2d} \prod_{i=1}^n \beta_i^{\alpha_i}
\]
Taking the $2d$-th root and rearranging gives
\[
   \mu > 2d \left( \prod_{i=1}^n \left( \frac{\beta_i}{\alpha_i} \right)^{\alpha_i} \right)^{1/2d}
\]
Since we assumed $f(x) \geq 0$ for all $x \in \mathbb{R}_+^n$, we can show a contradiction by constructing a specific
vector $x$ for which $f(x) < 0$. We choose the vector $x$ that corresponds to the equality condition of the AM-GM inequality
used above. This occurs when all the terms $y_i = \frac{\alpha_i}{\beta_i} x_i^{2d}$ are equal.
Let $x_i = \left( \frac{\beta_i}{\alpha_i} \right)^{1/2d}$.

Now we evaluate $f(x)$ for this specific $x$:
\begin{itemize}
   \item The sum term is
   \[
      \sum_{i=1}^n \beta_i x_i^{2d} = \sum_{i=1}^n \beta_i \left( \frac{\beta_i}{\alpha_i} \right) = \sum_{i=1}^n \alpha_i = 2d
   \]
   \item The monomial term is
   \[
      x^{\alpha} = \prod_{i=1}^n x_i^{\alpha_i} = \prod_{i=1}^n \left( \frac{\beta_i}{\alpha_i} \right)^{\alpha_i/2d} = \left( \prod_{i=1}^n \left( \frac{\beta_i}{\alpha_i} \right)^{\alpha_i} \right)^{1/2d}
   \]
\end{itemize}
Substituting these into $f(x)$:
\[
   f(x) = 2d - \mu \left( \prod_{i=1}^n \left( \frac{\alpha_i}{\beta_i} \right)^{\alpha_i} \right)^{1/2d}
\]
This implies that $2d - \mu \left( \prod_{i=1}^n \left( \frac{\alpha_i}{\beta_i} \right)^{\alpha_i} \right)^{1/2d} < 0$.
So, for our chosen $x$, we have found that $f(x) < 0$. This contradicts the premise that $f(x)$ is PSD. Therefore,
our assumption that the inequality does not hold must be false.
Both directions of the proof are complete. $\qed$
\end{proof}

\begin{theorem}
Let $\alpha(0), \alpha(1), \dots, \alpha(r) \in \mathbb{N}^n$ be exponent vectors. Let the exponent vector
$\beta \in \mathbb{N}^n$ be a convex combination of the $\alpha(j)$ vectors, such that $\beta = \sum_{j=0}^r \lambda_j \alpha(j)$,
with $\lambda_j > 0$ and $\sum_{j=0}^r \lambda_j = 1$. Consider the polynomial $f(X)$ with real coefficients $f_{\alpha} > 0$
defined as
\[
   f(X) = \sum_{j=0}^r f_{\alpha(j)} X^{\alpha(j)} - f_{\beta} X^{\beta}
\]
where $X = (X_1, \dots, X_n)$.

The polynomial $f(X)$ is positive semidefinite (PSD) for all $X \in \mathbb{R}_+^n$ if and only if
\[
   |f_{\beta}| \leq \prod_{j=0}^r \left( \frac{f_{\alpha(j)}}{\lambda_j} \right)^{\lambda_j}
\]
\end{theorem}

\begin{proof}[Proof Sketch]
The proof of the $(\Leftarrow)$ direction relies on the weighted Arithmetic Mean-Geometric Mean (AM-GM) inequality.
Define terms $Y_j$ and use the weights $\lambda_j$:
\[
   Y_j = \frac{f_{\alpha(j)}}{\lambda_j} X^{\alpha(j)}, \quad j = 0, \dots, r
\]
The weighted AM-GM inequality states that the arithmetic mean is greater than or equal to the geometric mean:
\[
   \sum_{j=0}^r \lambda_j Y_j \geq \prod_{j=0}^r Y_j^{\lambda_j}
\]
\begin{itemize}
   \item \textbf{Arithmetic Mean (LHS):}
   \[
      \sum_{j=0}^r \lambda_j Y_j = \sum_{j=0}^r f_{\alpha(j)} X^{\alpha(j)}
   \]
   \item \textbf{Geometric Mean (RHS):}
   \[
      \prod_{j=0}^r Y_j^{\lambda_j} = \prod_{j=0}^r \left( \frac{f_{\alpha(j)}}{\lambda_j} \right)^{\lambda_j} X^{\beta}
   \]
\end{itemize}
Combining these, the AM-GM inequality implies
\[
   \sum_{j=0}^r f_{\alpha(j)} X^{\alpha(j)} \geq \prod_{j=0}^r \left( \frac{f_{\alpha(j)}}{\lambda_j} \right)^{\lambda_j} X^{\beta}
\]
If we assume the theorem's condition holds ($f_{\beta} \leq \prod_{j=0}^r (\frac{f_{\alpha(j)}}{\lambda_j})^{\lambda_j}$),
then it follows that
\[
   \sum_{j=0}^r f_{\alpha(j)} X^{\alpha(j)} \geq f_{\beta} X^{\beta} \implies f(X) \geq 0
\]
This proves one direction of the theorem.
\end{proof}


\subsection*{Application to D-T Forms}
	\textbf{Problem.} Find the condition under which the function $f(x)$ is positive semidefinite for $x \in \mathbb{R}_+^n$:
\[
   f(x) = \sum_{i=1}^n \beta_i x_i^{2d} - \mu x^{\alpha}
\]
where $\beta_i, \mu > 0$ and the exponents satisfy $\sum_{i=1}^n \alpha_i = 2d$.

	\textbf{Solution via Theorem 2.2.}

We can map this problem to the general theorem's framework:
\begin{itemize}
   \item The "arithmetic mean" part is $\sum_{i=1}^n \beta_i x_i^{2d}$. The corresponding exponent vectors are $\alpha(i) = 2d e_i$, where $e_i$ is the $i$-th standard basis vector in $\mathbb{R}^n$. The coefficients are $f_{\alpha(i)} = \beta_i$.
   \item The "geometric mean" part is $\mu x^{\alpha}$. The exponent vector is $\beta = (\alpha_1, \alpha_2, \dots, \alpha_n)$.
   The coefficient is $f_{\beta} = \mu$.
   \item We must express $\beta$ as a convex combination of the $\alpha(i)$ vectors: $\beta = \sum_{i=1}^n \lambda_i \alpha(i)$.
\end{itemize}
This gives
\[
   (\alpha_1, \dots, \alpha_n) = \sum_{i=1}^n \lambda_i (0, \dots, 2d, \dots, 0) = (2d \lambda_1, \dots, 2d \lambda_n)
\]
which implies that the weights must be $\lambda_i = \frac{\alpha_i}{2d}$.
The condition $\sum \lambda_i = 1$ becomes $\sum \frac{\alpha_i}{2d} = 1$, which is exactly the given condition
$\sum \alpha_i = 2d$.

Now, we apply the inequality from the general theorem:
$$f_{\beta}\le\prod_{i=1}^{n}(\frac{f_{\alpha(i)}}{\lambda_{i}})^{\lambda_{i}}.$$
Substituting our specific terms:
$$\mu\le\prod_{i=1}^{n}(\frac{\beta_{i}}{\alpha_{i}/2d})^{\alpha_{i}/2d}=\prod_{i=1}^{n}(\frac{2d\cdot\beta_{i}}{\alpha_{i}} )^{\alpha_{i} / 2d}$$
To see that this is the same condition, we can raise both sides to the power of $2d$:
$$\mu^{2d}\le(\prod_{i=1}^{n}(\frac{2d\cdot\beta_{i}}{\alpha_{i}})^{\alpha_{i}/2d} )^{2d}=\prod_{i=1}^{n}(\frac{2d\cdot\beta_{i}}{\alpha_{i}})^{\alpha_{i}},$$ Since $\sum\alpha_i=2d$, this simplifies to:
$$\mu^{2d} \le \frac{(2d)^{2d} \prod \beta_i^{\alpha_i}}{\prod \alpha_i^{\alpha_i}}.$$
Thus
$$\mu^{2d}\prod \alpha_i^{\alpha_i}\le(2d)^{2d} \prod \beta_i^{\alpha_i}.\square$$

\section{Investigation of $M_{2d, p}(\underline{X}, \alpha)$ and Connection to SOBS and SOS}
\label{s03}
~~
\begin{proposition}
$M_{2d,0}(X, \alpha)$ is PSD D-T and SOBS.
\end{proposition}

\begin{proof}
See \cite{GM01}. (A simpler proof may be possible.)
% TODO: Add the proof for SOBS.
\end{proof}

\begin{theorem}
$M_{2d,p}(X, \alpha)$ with $d=1,2,3$, $p|2d$ is SOS for all $n \in \mathbb{N}$ and $\alpha \in \mathbb{N}^n$ with
$|\alpha|=2d$.
\end{theorem}

\begin{proof}
	\textbf{Case $d=1$:}
$M_{2,0}$ and $M_{2,1}$ are PSD quadratic forms and by Hilbert's theorem, they are SOS.

	\textbf{Case $d=2$:}
For $p=0$, these are PSD and SOBS by Proposition 4.1.
For $p=1$, we show that $M_{4,1}(X, \alpha)$ is a sum of squares (SOS).

For any number of variables $n$ and any nonnegative weights $\alpha = (\alpha_1, \dots, \alpha_n)$ such that $|\alpha| = 4$,
the form
\[
   M_{4,1}(X, \alpha) = \frac{\sum_{i=1}^n \alpha_i X_i^4}{4} - \left( \frac{\sum_{i=1}^n \alpha_i X_i}{4} \right)^4
\]
is SOS. We examine each case by the number of variables $n$:
\begin{enumerate}
   \item \textbf{Case $n=2$ (Binary Forms):} For $n=2$, $M_{4,1}$ is a PSD binary quartic form. By Hilbert's Theorem (1888),
   any nonnegative polynomial in two variables is a sum of squares. Therefore, $M_{4,1}$ is SOS for $n=2$.
   \item \textbf{Case $n=3$ (Ternary Quartics):} For $n=3$, $M_{4,1}$ is a PSD ternary quartic. Again, by Hilbert's Theorem,
   any nonnegative ternary quartic is a sum of squares. Therefore, $M_{4,1}$ is SOS for $n=3$.
   \item \textbf{Case $n=4$:}
      \begin{itemize}
         \item For uniform weights, $\alpha = (1,1,1,1)$:
         \[
            M_{4,1}(X, (1,1,1,1)) = \frac{\sum_{i=1}^4 X_i^4}{4} - \left( \frac{\sum_{i=1}^4 X_i}{4} \right)^4
         \]
         This specific form is known to be a sum of squares, verified computationally by semidefinite programming (SDP).
         % TODO: Find an explicit SOS decomposition for this form.
         \item For other $\alpha$ with $|\alpha|=4$ (e.g., $(2,1,1,0)$, $(3,1,0,0)$, $(4,0,0,0)$), all involve at most
         three non-zero entries, so they reduce to the $n=3$ case and are SOS.
      \end{itemize}
   \item \textbf{Case $n \geq 5$:}
      For any $\alpha = (\alpha_1, \dots, \alpha_n)$ with $|\alpha|=4$, there are at most four non-zero entries. Thus,
      $M_{4,1}(X, \alpha)$ only involves variables corresponding to non-zero weights, reducing to a case with at most $4$
      variables ($n \leq 4$). Since all cases for $n \leq 4$ are SOS, the proof is complete for $p=1$.
   \item For $p=2$, $M_{4,2}(X, \alpha)$ is SOS. For any $n$ and $\alpha$ with $|\alpha|=4$,
      \[
         M_{4,2}(X, \alpha) = \frac{\sum_{i=1}^n \alpha_i X_i^4}{4} - \left( \frac{\sum_{i=1}^n \alpha_i X_i^2}{4} \right)^2
      \]
      \begin{itemize}
         \item For $n=2$ and $n=3$, by Hilbert's Theorem, these are SOS.
         \item For $n=4$, with $\alpha = (1,1,1,1)$:
         \[
            M_{4,2}(X, (1,1,1,1)) = \frac{\sum_{i=1}^4 X_i^4}{4} - \left( \frac{\sum_{i=1}^4 X_i^2}{4} \right)^2
         \]
         This form is nonnegative and computation shows it is SOS.
         % TODO: Find an explicit SOS decomposition for this form.
         \item For other $\alpha$ with $|\alpha|=4$, all involve at most three non-zero entries, so they reduce to the
         $n<4$ case.
         \item For $n \geq 5$, any $\alpha$ has at most four non-zero entries, so the problem reduces to $n \leq 4$.
      \end{itemize}
\end{enumerate}

	\textbf{Case $d=3$:}
Consider the family of forms
\[
   M_{6,p}(X, \alpha)
\]
where $p \in \{0,1,2,3\}$, $|\alpha|=6$. We examine by the number of variables $n$:
\begin{enumerate}
   \item For $n=2$ (Binary Forms): By Hilbert's Theorem, all $M_{6,p}$ are SOS for $n=2$.
   \item For $n=3$ (Ternary Forms):
      \begin{itemize}
         \item For $p=0$, $M_{6,0}$ is known to be SOS.
         \item For $p=1$, analyze
         \[
            M_{6,1}(X, \alpha) = \frac{\sum_{i=1}^3 \alpha_i X_i^6}{6} - \left( \frac{\sum_{i=1}^3 \alpha_i X_i}{6} \right)^6
         \]
         Possible partitions of $6$ into $3$ parts for $\alpha \in \mathbb{N}^3$ are $(2,2,2)$, $(3,2,1)$, $(4,1,1)$, $(5,1,0)$,
         $(6,0,0)$.
         \item For $\alpha=(2,2,2)$: $M_{6,1}(X,(2,2,2))$ is known to be SOS.
         % TODO: Find an explicit SOS decomposition.
         \item For $\alpha=(3,2,1)$ and $(4,1,1)$: These can be reduced to $(2,2,2)$ by change of variables, so are SOS.
      \end{itemize}
   \item For $n=4$: Possible partitions of $6$ into $4$ parts for $\alpha$ are $(2,2,1,1)$, $(3,1,1,1)$, $(4,1,1,0)$, $(4,2,0,0)$, $(5,1,0,0)$, $(3,2,1,0)$, $(6,0,0,0)$. Cases with zeros reduce to $n \leq 3$.
      \begin{itemize}
         \item For $\alpha=(2,2,1,1)$ and $(3,1,1,1)$: These are known nonnegative forms and SDP decomposition shows
         they are SOS.
      \end{itemize}
   \item For $n=5$ and $n=6$: Partitions with zeros reduce to fewer variables. For $n=6$, only $(1,1,1,1,1,1)$ is without zeros.
\end{enumerate}
This completes the proof for Theorem 4.1.
\end{proof}

\begin{remark}
\begin{enumerate}
   \item It is of interest to investigate if Theorem 4.1 generalizes to any $d$.
   \item It would be interesting to know if these forms are in the SONC cone.
\end{enumerate}
\end{remark}

\section{Beyond the SOSONC Cone}
\label{s04}~~
\subsection{Specific Examples to Analyze}
\begin{itemize}
    \item $Q(X,Y,Z,W) = X^4 + Y^4 + Z^4 + W^4 - 4XYZW$
\end{itemize}

Similar as the Motzkin form, one can easily verify that both Choi Lam forms $Q$, $S$, and the modified Motzkin $M$ are indeed
PSD circuits (page 31, M. Schick thesis).
\[
    R(X,Y,Z) := X^6 + Y^6 + Z^6 - (X^4 Y^2 + X^4 Z^2 + X^2 Y^4 + X^2 Z^4 + Y^4 Z^2 + Y^2 Z^4)
    + 3X^2 Y^2 Z^2
\]
\[
    \hat{R}(X,Y,Z,W) := X^2(X-W)^2 + Y^2(Y-W)^2 + Z^2(Z-W)^2 + 2XYZ(X+Y+Z-2W)
\]
\[
    P(X,Y,Z) := 200 \left[ (X^3-4XZ^2)^2 + (Y^3-4YZ^2)^2 \right] + (Y^2-X^2)X(X+2Z)\left[ X(X-2Z)+2(Y^2-4Z^2) \right]
\]

Remarkably, there is the following relation between $Q$ and $R$ \cite[see page 5]{Rez07}:
\[
    \hat{R}(X-W, Y-W, Z-W, X+Y+Z-W) = 2Q(X,Y,Z,W)
\]
Since the PSD and SOS cones are closed under linear transformations, the identity above shows that $R$ is PSD (SOS)
if and only if $Q$ is PSD (SOS). Now, using the Newton polytope method one can prove that
$Q \in P^h_{4,4} \setminus \Sigma^h_{4,4}$. In this sense, the Robinson form $R$ admits a proof of
$\hat{R} \in P^h_{4,4} \setminus \Sigma^h_{4,4}$ following the Newton polytope method. (page 38)

\subsection{Separating Forms via Linear Transformation of the Variables}

\section{Positivestellensatz and relation to the cone of $2d$ powers}
\label{s05}
~~
\section{Application to polynomial optimization}
\label{s06}
~~
\begin{thebibliography}{99}

\bibitem{JBL01} J. B. Lasserre. \emph{Global optimization with polynomials and the problem of moments.} SIAM J. Optim., 11(3):796--817, 2001.

\bibitem{GM01} M. Ghasemi and M. Marshall, \emph{Lower bounds for polynomials using geometric programming}, SIAM J. Optim. 22 (2012), no. 2, 460--473.

\bibitem{GLM} M. Ghasemi, J.B. Lasserre, and M. Marshall, \emph{Lower bounds on the global minimum of a polynomial}, Comput. Optim. Appl. 57 (2014), no. 2, 387--402.

\bibitem{GM02} M. Ghasemi and M. Marshall, \emph{Lower bounds for a polynomial on a basic closed semialgebraic set using geometric programming}, 2013, Preprint, arxiv:1311.3726.

\bibitem{GM03} M. Ghasemi, M. Marshall, \emph{Lower Bound for a Polynomial on a product of hyperellipsoids using geometric programming} 2015, Preprint, arxiv:2510.03105.

\bibitem{DIdW01} M. Dressler, S. Iliman, and T. de Wolff, An Approach to Constrained Polynomial Optimization via Nonnegative Circuit Polynomials and Geometric Programming, \emph{Journal of Symbolic Computation} (2019) \textbf{91}, 149-172

\bibitem{DIdW02} M. Dressler, S. Iliman, and T. de Wolff. \emph{A Positivstellensatz for Sums of Nonnegative Circuit Polynomials}. SIAM J. Appl. Algebra Geom., 1(1):536--555, 2017.

\bibitem{IdW} Sadik Iliman and Timo De Wolff, \emph{Lower Bounds for Polynomials with Simplex Newton Polytopes Based on Geometric Programming,} SIAM Journal on Optimization 26, no. 2 (January 2016): 1128--46, \url{https://doi.org/10.1137/140962425}

\bibitem{MSc01} Moritz Schick \emph{Sums of Squares plus Sums of Nonnegative Circuit Polynomials}, PhD. Dissertation. 2025.

\bibitem{Rez07} B. Reznick. \emph{On Hilbert's construction of positive polynomials}, 2007. Preprint, arXiv:0707.2156.

\bibitem{DresslerSchick2025} M. Dressler, S. Kuhlmann and Moritz Schick, \emph{Study of the cone of sums of squares plus sums of nonnegative circuit forms} Advances in Geometry, 25(1), 127--146 (2025)

\end{thebibliography}

\bigskip

\noindent\textbf{Mehdi Ghasemi:} mehdi.ghasemi@usask.ca\\
Edmonton Police Service

\noindent\textbf{Salma Kuhlmann:} salma.kuhlmann@uni-konstanz.de\\
University of Konstanz

\end{document}
